\emph{Setting.} The labels are those of Theorem \claimref{synthesis.wick-rotated-boost-selects-observer-tube-point}, and the de Sitter gravitational field is that of Theorem \claimref{synthesis.de-sitter-horizon-translation-is-conformal-null-rotation}, rewritten in these labels: $e_{1}$ there is the time direction ($e_{0}$ in that theorem), $e_{+}$ plays the role of $\nu$, and $\ip{\nO}{\ninf} = 1$, $Q(\nO) = Q(\ninf) = 0$. A constant tractor $I$ gives the scale $\sigma(x) = \ip{X(x)}{I}$; for $I_{\Lambda} = \ninf + \tfrac{\Lambda}{6}\nO$, $\ip{X(x)}{I_{\Lambda}} = 1 - \tfrac{\Lambda}{12}Q(x)$ because $\ip{X(x)}{\ninf} = 1$ and $\ip{X(x)}{\nO} = -\tfrac12 Q(x)$. A bivector $B$ generates a symmetry of the gravitational field $\hb = \sigma$ exactly when it commutes with $I_{\Lambda}$, since it then fixes $\sigma$ (that theorem, (i)), and it fixes a compatible tube point $m$ exactly when the infinitesimal action $[B, X(m)]$ is a complex multiple of $X(m)$, the complex null line of $X(m)$ determining $m$ (Proposition \claimref{one-relation.spacetime-manifold-vacuum-lines}, extended complex-linearly as in the cited tube-point theorem). We take $K = e_{1}e_{2}$, twice the stored generator; its velocity field at $x$ is $K \cdot x$, which has unit rapidity per unit parameter. For $m = x_{\perp} + I_{7}\rho u$, $Q(m) = Q(x_{\perp}) + (I_{7}\rho)^{2}Q(u) = -|x_{\perp}|^{2} - \rho^{2}$ (as $x_{\perp} \cdot u = 0$), so
\[
\operatorname{Re}X(m) = \nO + x_{\perp} + \tfrac12(|x_{\perp}|^{2} + \rho^{2})\ninf, \qquad \operatorname{Im}X(m) = \rho u .
\]

\emph{(i)} $K$ and $I_{\Lambda}$ are a bivector of $\operatorname{span}(e_{1}, e_{2})$ and a vector orthogonal to that plane. A bivector $a \wedge b$ with $a, b$ orthogonal to $I_{\Lambda}$ commutes with $I_{\Lambda}$, and one built from $\operatorname{span}(e_{1}, e_{2})$ or from its orthogonal complement commutes with $K$. The orthogonal complement of $\operatorname{span}(e_{1}, e_{2}, I_{\Lambda})$ is $\operatorname{span}(e_{3}, e_{4}, W)$ with $W = \ninf - \tfrac{\Lambda}{6}\nO$: indeed $\ip{W}{I_{\Lambda}} = \tfrac{\Lambda}{6} - \tfrac{\Lambda}{6} = 0$ and $Q(W) = -2\cdot\tfrac{\Lambda}{6} = -\tfrac{\Lambda}{3}$ (for $\Lambda = 0$ use $\ninf$ in place of $W$; $I_{0} = \ninf$ and the complement of $\operatorname{span}(e_{1}, e_{2})$ in $\ninf^{\perp}$ is $\operatorname{span}(e_{3}, e_{4}, \ninf)$, degenerate). So $K$, $e_{3}e_{4}$, $e_{3}W$, $e_{4}W$ commute with both. Nothing else does: a bivector commuting with $K$ has no mixed part $e_{a} \wedge c$ with $a \in \{1, 2\}$ and $c \in V := \operatorname{span}(e_{3}, e_{4}, e_{+}, e_{-})$ (such a part anticommutes with $e_{1}e_{2}$), so it is a multiple of $K$ plus a bivector of $V$; and the bivectors of the nondegenerate four-dimensional $V$ fixing the nonzero vector $I_{\Lambda} \in V$ form a three-dimensional algebra, whether $I_{\Lambda}$ is null or not (the orbit of a nonzero vector under $\mathrm{SO}(V)$ is three-dimensional). The check confirms the dimension $4$ of the kernel of $B \mapsto ([B, I_{\Lambda}], [B, K])$ on the fifteen bivectors exactly for $\Lambda \in \{12, 3, 48, \tfrac43, 0, -12, -\tfrac34\}$. For orthogonal $e_{3}, e_{4}, W$, $e_{3}We_{4}W = -Q(W)e_{3}e_{4}$ and $e_{4}We_{3}W = Q(W)e_{3}e_{4}$, so $[e_{3}W, e_{4}W] = -2Q(W)\,e_{3}e_{4}$, and cyclically: the three transverse generators close on $\so(\operatorname{span}(e_{3}, e_{4}, W))$, which is compact exactly when $Q(W) < 0$, that is $\Lambda > 0$ (all three directions of square $-1$, an $\so(3)$), is $\so(1,2)$ for $\Lambda < 0$, and for $\Lambda = 0$ is the algebra of the rotation $e_{3}e_{4}$ and the translations $P_{e_{3}}, P_{e_{4}}$ ($P_{a} = \tfrac12 a\ninf$) of the bifurcation plane.

\emph{(ii)} Let $\Lambda > 0$ and $m$ compatible. The three transverse generators $B$ commute with $e_{1}$ and $e_{2}$, so $[B, \operatorname{Im}X(m)] = 0$, and $[B, \operatorname{Re}X(m)]$ is a real vector. If $[B, X(m)] = cX(m)$ with $c = \alpha + I_{7}\beta$, comparing imaginary parts gives $0 = \beta\operatorname{Re}X + \alpha\operatorname{Im}X$, so $\alpha = \beta = 0$ because $\operatorname{Re}X$ has an $\nO$-component and $\operatorname{Im}X$ does not; hence $B$ fixes $m$ if and only if $[B, \operatorname{Re}X(m)] = 0$. The vectors fixed by $\so(\operatorname{span}(e_{3}, e_{4}, W))$ are those of $\operatorname{span}(e_{1}, e_{2}, I_{\Lambda})$, and $\operatorname{Re}X(m)$ has no component along $e_{1}$, $e_{2}$. So all three fix $m$ if and only if $\operatorname{Re}X(m)$ is a multiple of $I_{\Lambda}$. Writing $\nO = \tfrac{3}{\Lambda}(I_{\Lambda} - W)$ and $\ninf = \tfrac12(I_{\Lambda} + W)$, the components of $\operatorname{Re}X(m)$ along $e_{3}, e_{4}$ are those of $x_{\perp}$ and along $W$ it is $-\tfrac{3}{\Lambda} + \tfrac14(|x_{\perp}|^{2} + \rho^{2})$; both vanish if and only if $x_{\perp} = 0$ and $\rho^{2} = 12/\Lambda$. Then $\operatorname{Re}X(m) = \nO + \tfrac{6}{\Lambda}\ninf = \tfrac{6}{\Lambda}I_{\Lambda} = \tfrac{\rho_{\Lambda}^{2}}{2}I_{\Lambda}$, the converse of the boost-fixed vector computed in Theorem \claimref{synthesis.positive-structure-and-natural-scale-actions-do-not-select-the-almost-einstein-scale}(i). The boost $K$ changes $\eta$ and fixes $x_{\perp}$ and $\rho$, and it moves every compatible tube point ($[K, X(m)]$ is not a multiple of $X(m)$), so the fixed set is one boost orbit, $\{I_{7}\rho_{\Lambda}u\}$, whose observers $x(m) = \rho_{\Lambda}(\sinh\eta\,e_{1} + \cosh\eta\,e_{2})$ (that theorem, (iii)) fill the hyperbola $Q(x) = -12/\Lambda$ of the right wedge. The check confirms the fixed set on a grid of $36$ compatible points (three feet, six scales, three rapidities given by rational half-rapidity parameters) for each of four values of $\Lambda$, finding the predicted point fixed and every other point moved by some transverse generator.

\emph{(iii)} At $x = x(m)$ with $x_{\perp} = 0$ and $\eta = 0$, $x = \rho e_{2}$ and $K \cdot x = -\rho e_{1}$, with $Q = \rho^{2}$; by boost invariance the same holds along the orbit. With $\hb = \sigma$, the physical vector of the background velocity $a$ is $\hb^{-1}(a) = a/\sigma$ (Fact \claimref{synthesis.imported-gtg-dictionary}(i)), so the physical norm is $N(\rho) = \rho/\sigma(\rho e_{2}) = \rho/(1 + \tfrac{\Lambda}{12}\rho^{2})$, using $Q(\rho e_{2}) = -\rho^{2}$. For $\Lambda = 3H^{2} > 0$, $\sigma - H\rho = 1 + \tfrac{H^{2}\rho^{2}}{4} - H\rho = (1 - \tfrac{H\rho}{2})^{2}$, so
\[
1 - H N(\rho) = \frac{(1 - H\rho/2)^{2}}{\sigma} \ge 0 ,
\]
with equality if and only if $\rho = 2/H = \rho_{\Lambda}$. By Theorem \claimref{synthesis.de-sitter-static-boost-frozen-point}(iii) the surface gravity of the de Sitter field, with the static flow normalised to unit norm at its observer, is $\kappa_{\Lambda} = H$; so $N \le 1/\kappa_{\Lambda}$ with equality exactly on the orbit of (ii), where $\sigma = 2$. At the chart distance $\rho = 1/H$, $\sigma = \tfrac54$ and $N = \tfrac{4}{5H}$. For $\Lambda = 0$, $N = \rho$; for $\Lambda < 0$, $N = \rho/(1 - \tfrac{|\Lambda|}{12}\rho^{2})$ on $\sigma > 0$; both are strictly increasing, with no maximum.

\emph{(iv)} Parametrise the orbit as $x(\lambda) = \rho w(\lambda)$, $w = \sinh\lambda\,e_{1} + \cosh\lambda\,e_{2}$, $u = dw/d\lambda$, $du/d\lambda = w$. The unit physical velocity is $v = u$: $x' = dx/d\tau = \hb(v) = \sigma v$ with $d\tau/d\lambda = \rho/\sigma = N$, and $Q(u) = 1$. By Fact \claimref{synthesis.imported-gtg-dictionary}(iii) the orbit is a free-fall worldline exactly when $A := dv/d\tau + \Omega(x') \cdot v = 0$. Here $dv/d\tau = \tfrac{\sigma}{\rho}w$; with $\Omega(a) = -(a \wedge \nabla\sigma)/\sigma$ and $\nabla\sigma = -\tfrac{\Lambda}{6}x$ (the vector derivative of $Q(x)$ being $2x$), $\Omega(x') = -(v \wedge \nabla\sigma)$ and $\Omega(x') \cdot v = -v\,(\nabla\sigma \cdot v) + \nabla\sigma\,Q(v) = \nabla\sigma = -\tfrac{\Lambda\rho}{6}w$, since $w \cdot u = 0$. So
\[
A = \Bigl(\frac{1 + \Lambda\rho^{2}/12}{\rho} - \frac{\Lambda\rho}{6}\Bigr)w = \frac{1 - \Lambda\rho^{2}/12}{\rho}\,w ,
\]
zero if and only if $\Lambda\rho^{2} = 12$, never for $\Lambda \le 0$. The check evaluates $A$ in $\Cl(2,4)$ from the stated $\Omega$ and the bivector-vector product at three rapidities, six scales and six values of $\Lambda$. The rotor $R = 1 + e_{2}W$ at $\Lambda = 3$ ($Q(W) = -1$, $R\tilde R = 2$) fixes $I_{3}$ (it is built from vectors orthogonal to $I_{3}$), carries $K$ to $e_{1}W$, the static boost $B_{\Lambda} = e_{1} \wedge W$ of Theorem \claimref{synthesis.de-sitter-static-boost-frozen-point} in these labels, and carries the stored static observer's worldline $t \mapsto te_{1}$ onto the $K$-orbit $Q(x) = -4 = -12/\Lambda$ (checked at four values of $t$, the images normalised by $\ip{\cdot}{\ninf} = 1$): the orbit of (ii) is the stored static observer seen through a symmetry of the field.

\emph{(v)} On $\operatorname{span}(e_{3}, e_{4})$, with chart coordinates $(y, z)$ and $r^{2} = y^{2} + z^{2} = -Q$, $\sigma = 1 + \tfrac{\Lambda}{12}r^{2}$ and the physical metric is $\sigma^{-2}(dy^{2} + dz^{2})$. For a metric $e^{2f}(dy^{2} + dz^{2})$ the Gauss curvature is $-e^{-2f}\Delta f$; with $f = -\ln\sigma$ this is $\sigma^{2}\Delta\ln\sigma = \sigma\Delta\sigma - |\nabla\sigma|^{2} = \tfrac{\Lambda}{3}(1 + \tfrac{\Lambda}{12}r^{2}) - \tfrac{\Lambda^{2}}{36}r^{2} = \tfrac{\Lambda}{3}$. So for $\Lambda > 0$ the bifurcation plane, completed by its point at infinity, is the round sphere of radius $\ell = (3/\Lambda)^{1/2} = 1/\kappa_{\Lambda}$; its rotation algebra is the $\so(3)$ of (i), restricted (the generators fix $\sigma$, hence act by isometries). The circle of chart radius $r$ about the origin has physical radius $r/\sigma = N(r)$, the function of (iii), so the largest circle (the equator, of physical radius $\ell$) is at chart radius $\rho_{\Lambda}$. The check computes $\sigma\Delta\sigma - |\nabla\sigma|^{2}$ by central differences, exact on the quadratic $\sigma$, at four points for five values of $\Lambda$.

\emph{(vi)} For $\Lambda \le 0$ the argument of (ii) applies with $\so(1,2)$ or the transverse translations: the transverse generators again commute with $e_{1}, e_{2}$, so a fixed $m$ needs $[B, \operatorname{Re}X(m)] = 0$ for all three; for $\Lambda < 0$ the fixed vectors are again $\operatorname{span}(e_{1}, e_{2}, I_{\Lambda})$, and the $\nO$ and $\ninf$ components of $\operatorname{Re}X(m)$ are $1$ and $\tfrac12(|x_{\perp}|^{2} + \rho^{2}) > 0$, in the ratio $\tfrac{\Lambda}{6} < 0$ only if $\rho^{2} < 0$; for $\Lambda = 0$, $[P_{e_{3}}, \operatorname{Re}X(m)]$ has the component $e_{3}$ times a nonzero multiple of $\ip{\operatorname{Re}X}{\ninf} = 1$, so the translation moves every point. No compatible tube point is fixed. (iii) and (iv) gave the rest.

\emph{(vii)} $\ip{u}{I_{\Lambda}} = 0$, so $\sigma(m) = \ip{\operatorname{Re}X(m)}{I_{\Lambda}} = 1 + \tfrac{\Lambda}{12}(|x_{\perp}|^{2} + \rho^{2})$ is real and positive for $\Lambda > 0$. Since $Q(X(m)) = 0$ (complex bilinear $Q$) and $\ip{X(m)}{I_{\Lambda}} = \sigma(m)$, $V(m) = X(m)/\sigma(m) - \tfrac{3}{\Lambda}I_{\Lambda}$ has $\ip{V}{I_{\Lambda}} = 1 - 1 = 0$ and $Q(V) = -\tfrac{6}{\Lambda} + \tfrac{9}{\Lambda^{2}}\cdot\tfrac{\Lambda}{3} = -\tfrac{3}{\Lambda}$, the normalisation of the Stelle-West vector in Theorem \claimref{dark-energy.frozen-point-makes-stelle-west-vector}(ii), here at a complex point. Its real and imaginary parts are $\operatorname{Re}X/\sigma - \tfrac{3}{\Lambda}I_{\Lambda}$ and $\rho u/\sigma$; the imaginary part of $Q(V)$ is $2\ip{\operatorname{Re}V}{\operatorname{Im}V} = 0$ and its real part gives $Q(\operatorname{Re}V) - Q(\operatorname{Im}V) = -3/\Lambda$, with $Q(\operatorname{Im}V) = \rho^{2}/\sigma^{2}$. $\operatorname{Re}V$ is orthogonal to $I_{\Lambda}$ and lies in $\operatorname{span}(\nO, \ninf, e_{3}, e_{4})$, hence in $\operatorname{span}(W, e_{3}, e_{4})$, negative definite for $\Lambda > 0$; so $\rho^{2}/\sigma^{2} = \tfrac{3}{\Lambda} + Q(\operatorname{Re}V) \le \ell^{2}$, with equality exactly when $\operatorname{Re}V = 0$, that is when $\operatorname{Re}X(m)$ is a multiple of $I_{\Lambda}$, the condition of (ii). This extends (iii) to every foot. The conjugate $\bar m$ has the same real part, so by the argument of (ii) it is fixed by the same symmetries. The check verifies every identity on the grid of (ii), with the conjugate points.

\emph{(viii)} $E = e_{+}e_{-}$ has $E^{2} = 1$, and $R = 1 + tE$ ($|t| < 1$) has $R\tilde R = 1 - t^{2}$; it commutes with $e_{1}, \dots, e_{4}$ and acts on $\operatorname{span}(\nO, \ninf)$ by $\nO \mapsto s\nO$, $\ninf \mapsto s^{-1}\ninf$ for some $s > 0$ ($s$ is $\ip{R\nO\tilde R}{\ninf}/(1 - t^{2})$). Then $I_{\Lambda} = \ninf + \tfrac{\Lambda}{6}\nO \mapsto s^{-1}(\ninf + \tfrac{s^{2}\Lambda}{6}\nO) = s^{-1}I_{s^{2}\Lambda}$, and $\operatorname{Re}X(I_{7}\rho e_{1}) = \nO + \tfrac{\rho^{2}}{2}\ninf \mapsto s(\nO + \tfrac{(\rho/s)^{2}}{2}\ninf)$, $\operatorname{Im}X \mapsto s\cdot(\rho/s)e_{1}$: the tube point $I_{7}\rho e_{1}$ goes to $I_{7}(\rho/s)e_{1}$, and $(\rho/s)^{2}(s^{2}\Lambda) = \rho^{2}\Lambda$. At $\Lambda = 12$, $I_{12} = \ninf + 2\nO = 2e_{+}$ and $\rho_{12} = 1$, $X(I_{7}e_{1}) = e_{+} + I_{7}e_{1}$: the fixed tube point is $I_{7}e_{1}$, the point at which the stored positive structure is the twistor form polarised (Theorem \claimref{synthesis.positive-structure-is-twistor-form-at-imaginary-time}). The check verifies these at three values of $t$ and four of $\Lambda$.

\emph{The consequence.} The open problem lists two candidates for the scale and one possible source. In the de Sitter field: the frozen-point candidate $\rho^{2} = 12/\Lambda$ is the tube point fixed by the field's symmetries commuting with the boost, (ii); the surface-gravity candidate, read as the requirement that the boost generator normalised by $\kappa_{\Lambda}$ have unit physical norm at the observer, holds exactly there, (iii), a double root; and the round metric of the bifurcation surface has its rotation algebra and equator there, (v). In chart length the factor $2$ between $\rho_{\Lambda}$ and $1/\kappa_{\Lambda}$ remains, $\rho_{\Lambda} = 2/\kappa_{\Lambda}$, and it is the local scale $\sigma(x(m)) = 2$; equivalently $\rho_{\Lambda} = \sigma(m)\,|\operatorname{Im}V(m)|$ with $\sigma(m) = 2$ by (vii). Read as a chart length, $\rho = 1/\kappa_{\Lambda}$ gives $N = \tfrac45\ell$, an orbit singled out by none of the three. With no field ($\Lambda = 0$), or $\Lambda < 0$, none of the three selects a scale, and the dilations of the cited tube-point theorem leave $\rho$ free. \emph{Scope.} The coincidence is one of a field with one horizon. For the Schwarzschild-de Sitter lapse $f(r) = 1 - 2M/r - \tfrac{\Lambda}{3}r^{2}$ with $\Lambda = \tfrac37$, $M = \tfrac37$, $rf(r) = -\tfrac17(r - 1)(r - 2)(r + 3)$, so the horizons are at $r_{b} = 1$, $r_{c} = 2$, and $\kappa = \tfrac12|f'(r_{h})|$ gives $\kappa_{b} = \tfrac27$, $\kappa_{c} = \tfrac{5}{28}$, $\kappa_{b}/\kappa_{c} = \tfrac85$ (a ratio independent of the normalisation of the Killing field); one static region with one Killing flow has two surface gravities, so the surface-gravity candidate depends on the horizon and cannot agree with a single frozen-point scale at both. That field is not a pure dilation of the background, and nothing else here is claimed for it.

\emph{Reading (argued, not part of the statement).} Which premise fixes $\rho$: on this evidence, the gravitational field itself. The de Sitter-invariant Hadamard state is unique for a massive scalar (Allen's theorem, imported in the store as synthesis.imported-de-sitter-invariant-hadamard-state-unique; without the Hadamard condition the de Sitter-invariant states form the alpha family; the series' boundary modes are CAR fermions, outside the theorem's hypotheses, so this step is imported across a hypothesis gap), so its symmetry group contains the $\so(1,1) \oplus \so(3)$ of (i); any rule that assigns a tube point to a state and a wedge covariantly must then return a point fixed by the $\so(3)$, which by (ii) is the frozen-point tube point. The local KMS temperature of the boost-KMS state at rapidity period $2\pi$ on the orbit through $x(m)$ is $1/(2\pi N)$ (Tolman's relation, cited from memory), so by (iii) the selected observer is the coldest of the wedge, at the de Sitter temperature $H/2\pi$, and the free-fall one by (iv); its Euclidean boost orbit is a great circle of the Euclidean section, of length $2\pi\ell$. For the Schwarzschild field these tools do not apply: it is not a pure dilation of the background (Theorem \claimref{synthesis.de-sitter-horizon-translation-is-conformal-null-rotation}(vi)), and the static Killing norm normalised by $\kappa = 1/4M$ reaches $1$ only at infinity, so the analogue of (iii) selects the asymptotic observer, not a finite distance; not computed. The selection rule (covariance of the state's tube point under the field's symmetries commuting with the boost) is adopted as a costed clause, Premise \claimref{synthesis.tube-point-clause-invariant-under-boost-commuting-symmetries}, which carries the hypothesis gap just named and decides fields with one horizon only; (i) to (viii) do not use it. By (vii) it has an equivalent form in a datum the series already holds: the observer's tube point is the one at which the complexified Stelle-West vector of the MacDowell-Mansouri reduction, with $X(m)$ normalised by $\sigma(m)$, is purely imaginary, $V(m) = I_{7}\ell u$ (the tube point is not $V$: $\operatorname{Re}X(m) = \tfrac{6}{\Lambda}I_{\Lambda}$ and $\rho_{\Lambda} = \sigma(m)|\operatorname{Im}V| = 2\ell$); that condition is stated for every field with $\Lambda \neq 0$ at the level of the tractor fibre, but only for $\Lambda > 0$ does it have a solution, and only for the de Sitter field is it computed here. By (viii) the scale $\rho$ is the boost parameter of the ascent plane's null pair, the same kind of parameter as the rung unit of synthesis.rung-unit-transverse-to-norm one completion higher (argued analogy): what the de Sitter field fixes is the unit-independent product $\rho^{2}\Lambda = 12$, and in the balanced frame, where the frozen point is a multiple of one orthonormal vector of the ascent plane, the fixed point is the stored $I_{7}e_{1}$, as the balanced rung frame $\mu^{2} = \Lambda/3$ is the mirror-symmetric one there. So symmetry selection of this kind ties the unit to $\Lambda$, a Hubble-scale unit, and cannot supply an independent scale. The unit $I_{7}$ in $m = x + I_{7}y$ is the complex unit of the complexified vector space, the central pseudoscalar of $\Cl(3,4)$ (the store's convention for complexified points, Theorem \claimref{std-algebra.complexification-is-the-odd-extension}), not the conformal spinors' $I_{6}$; read with $I_{6}$, which is $\pm I_{7}$ on the two chirality halves, the forward tube of one half is the backward tube of the other, and by (vii) the selected point and its conjugate are fixed together, so the selection itself does not choose a chirality; the choice of which of the two is forward does.
